% TeX root=../main.tex

\begin{table}[!htb]
	\caption{Palatalizações do \glsxtrlong{eslavocomum} para o eslavo
		antigo}\label{tab:palatalizações}
	\begin{center}
		\begin{tabular}[c]{llll}
			\toprule
			                       & *tj            & \ocsex{št}                            & \cirilex{шт}                             \\
			Dentais                & *dj            & \ocsex{žd}                            & \cirilex{жд}                             \\
			                       & *sj            & \ocsex{s}̌                             & \cirilex{ш}                              \\
			\midrule
			Labiais                & *pj, *bj, *mj  & \ocsex{pl'}, \ocsex{bl’}, \ocsex{ml’} & \cirilex{пл}, \cirilex{бл}, \cirilex{мл} \\
			\midrule
			Soantes                & *nj, *rj, *lj  & \ocsex{n'}, \ocsex{r’}, \ocsex{l’}    & \cirilex{н},\cirilex{р}, \cirilex{л}     \\
			\midrule
			Velares:               & *kj, *kī̆, *kē̆  & \ocsex{C}                             & \cirilex{ч}                              \\
			primeira palatalização & *gj, *gī̆, *gē̆  & \ocsex{Z}                             & \cirilex{ж}                              \\
			\midrule
			Velares:               & *kě (*kai)     & \ocsex{c}                             & \cirilex{ц}                              \\
			segunda palatalização  & *\ocstrans{gE} & \ocsex{dz/z}                          & \cirilex{ꙁ}\slash{}\cirilex{з}           \\
			                       & *\ocsex{xE}    & \ocsex{s} (\ocsex{š} em morávio)      & \cirilex{с}                              \\
			\bottomrule
		\end{tabular}
	\end{center}
\end{table}
